% =====================================================================
%  Chapter 6 — Discussion  (~1,500 words)
%  DRAFT generated 2026-07-18.
% =====================================================================
\chapter{Discussion}
\label{ch:discussion}

\section{What the Findings Mean Together}
\label{sec:disc-synthesis}

Read separately, the results of Chapter~\ref{ch:eval} are a mixed report
card: the gated policy beats always-retrieve but trails closed-book on
multiple choice, leads on open-ended, and costs more wall-clock time than
either. Read together, they assemble into a single argument that revises
what a retrieval gate is for.

The starting observation is that \emph{retrieval is not a free safety net}.
On questions the model already answers correctly, injecting
topically-relevant-but-non-decisive context makes it worse at a ratio of
3.5 broken answers for every one fixed
(Section~\ref{sec:eval-distraction}). The mechanism is displacement, not
noise: the model defers to the provided context, as RAG prompting
instructs it to, and abandons parametric knowledge that was correct. The
standard RAG assumption, that context is at worst neutral, is false for this
model class.

A natural objection is that this ``distraction'' is really \emph{absence}:
perhaps retrieval hurts simply because the corpus does not contain the
answer. Three lines of evidence rule that out as the primary cause. If the
corpus lacked answers, retrieval could never turn a wrong answer right, yet
it \emph{fixes} \fixLlama{}, \fixQwenSeven{} and \fixQwenFourteen{}
closed-book errors at 3B, 7B and 14B. Retrieval measurably \emph{helps} the
weakest model on open-ended questions (P1 \fOnePoneOpen{} over P3
\fOnePthreeOpen{}), which is impossible if the information is not there. And
decisively, the harm \emph{grows} with capability
(Section~\ref{sec:scaling-mechanism}): a corpus that merely lacked
information would damage a stronger model \emph{less}, since that model can
fall back on what it knows, yet the strongest model is harmed the most,
because it has the most correct parametric knowledge to be displaced.
Absence cannot produce that gradient; displacement is its signature.

If retrieval can be poison, then the value of a gate is not primarily
\emph{efficiency}, the framing the adaptive-retrieval literature has
largely adopted, but \emph{protection}. The sharpest evidence is the
protective-skip result (Section~\ref{sec:eval-protective}): on precisely the
questions where retrieval breaks a correct answer, the gate's skip decision
scored 19/19, its retrieve decision 3/27. The gate's real job, on this
evidence, is to detect the boundary of the model's parametric knowledge and
defend everything inside it.

That reframing also resolves the apparent MCQ/open-ended contradiction
without special pleading. MIRAGE's multiple-choice questions sit largely
\emph{inside} the 3B model's parametric knowledge, so retrieval can only
distract and closed-book wins; PubMedQA's research questions sit largely
\emph{outside} it, so retrieval adds information and the gated policy wins.
One boundary, two sides, one gate whose task is to tell them apart. Even the
refuted ablation hypothesis (Section~\ref{sec:eval-ablation}) folds into the
same picture: probe-boosting variants fail because they retrieve more, and
on a benchmark where marginal retrievals hurt, the most conservative voting
rule wins. Three independent analyses, namely the paired study, the
protective skip and the ablation, converge on the same conclusion from different
directions, which is the kind of agreement that survives scrutiny better
than any single number.

\section{Limitations and Threats to Validity}
\label{sec:disc-limitations}

The following limitations are stated in full, ordered roughly by how much
they constrain the conclusions.

\paragraph{Statistical power.} $n=\nQuestions{}$ per condition, one seed, one
run per configuration. The headline accuracy gap of \diffPfivePone{}pp
carries a 95\% CI of \diffPfivePoneCI{}, which crosses zero: this study is
\emph{not powered to establish the accuracy claim}, and the abstract and
conclusion phrase it accordingly (``consistent with an improvement''). The
retrieval-cost halving is a count and is exact for this sample. The
high-risk stratum ($n=6$, Wilson CIs $\approx \pm 30$pp) supports no
between-policy comparison at all.

\paragraph{Scaling arm: three points, one family.} The scaling study
(Chapter~\ref{ch:scaling}) runs the full pipeline at 3B, 7B
and 14B and confirms calibration non-transfer prospectively at each new scale.
The ordering P3 $>$ P5 $>$ P1 reproduces on multiple-choice at all three scales,
and the open-ended ordering \emph{inverts} between the smallest model
(retrieval helps) and the two larger ones (retrieval hurts). But three points
are still a curve through one Llama model and two Qwen models, not a claim about
model families in general, and the two larger points come from the same family,
so a family effect cannot be separated from a scale effect between them. The
curve is also visibly \emph{saturating}, since the closed-book gain from 7B to
14B (\scaleGainLarge{}pp) is a fraction of the 3B-to-7B gain
(\scaleGainSmall{}pp), so extrapolating it to the higher safety bars would be
unwarranted.

\paragraph{Corpus completeness and retrieval quality.} The distraction result
does not exonerate the corpus. The index is \textasciitilde426K chunks drawn
from MedRAG textbooks and a 300K-passage PubMed slice; StatPearls, part of the
full MedRAG corpus, was unavailable at build time, so this is a subset rather
than the maximal medical corpus. The chunk-relevance analysis
(Section~\ref{sec:eval-distraction}) found the answer-bearing passage is
retrieved for only a minority of questions even when the gated answer is
correct, so there is genuine retrieval-quality headroom: a fuller corpus or a
stronger retriever could surface the decisive passage more often and reduce the
damage on the questions where retrieval currently breaks a correct answer.
The claim defended here is therefore the narrower one, that on \emph{this}
corpus and retriever, and for these models, the dominant failure mode is
displacement of correct parametric knowledge rather than absence of information
(the three-model evidence above), rather than that retrieval quality is
irrelevant.
Whether a substantially better retrieval stack would flip any of the MCQ
orderings is untested and is left to future work.

\paragraph{Retrieval budgets are not matched across models.} Llama's gated
policy retrieved on \retPfiveMcq{} of queries, Qwen's on
\retQwenPfiveMcq{}, a gap of \budgetGapMcq{}pp. Thresholds for Qwen were
refitted to reproduce the 3B's \emph{per-gate} budget and the two tunable
members landed close to target, but the hallucination probe is threshold-free
on multiple-choice (two drafts agree on a letter or they do not) and so
cannot be recalibrated at all. Its vote rate fell from \probeFireLlamaMcq{}
to \probeFireQwenMcq{} because the larger model self-agrees far more often,
and with one member of a majority-of-three effectively silent the ensemble
collapses to a conjunction of the other two and cannot reach the intended
rate. The consequence is precise and worth stating plainly: \emph{``selective
retrieval beats always-retrieve at a matched budget'' is established within
each model and not across them}, because the cross-model comparison confounds
scale with retrieval rate. Only the closed-book policies, which retrieve 0\%
by definition, are cleanly comparable between models.

\paragraph{Gate thresholds are task-specific, and were not treated as such.}
Thresholds were fitted on multiple-choice drafts and reused unchanged on
open-ended drafts. That is unsound: a multiple-choice draft is a single
letter and an open-ended draft is a paragraph, and their signal
distributions do not coincide. On Qwen the reuse was catastrophic: the
entropy signal's minimum (\qwenOpenEntropyMin{}) lies above its threshold
(\qwenOpenEntropyTau{}) and the margin signal's maximum
(\qwenOpenMarginMax{}) lies below its own (\qwenOpenMarginTau{}), so every
question retrieved and gated P5 degenerated into always-retrieve P1. No
open-ended selective-retrieval result can be drawn from that run. Llama
avoided total saturation only because its $\tau = 0.70$ happens to fall
inside its open-ended distribution, which is luck rather than design; its
open-ended retrieval rate of \retPfiveOpen{} still overshoots the intended
budget substantially. The general lesson is that a gate threshold is a
property of a model \emph{and} a task format, and the method as implemented
encodes only the former.

\paragraph{The open-ended metric barely discriminates.} Every open-ended
condition in this study scores between $0.19$ and $0.26$ mean token-$F_1$
while the underlying behaviour ranges from answering essentially every
question to declining \retPoneOpen{} of them. Declining costs almost
nothing: on always-retrieve, refusals average $0.217$ and substantive
answers $0.244$. Token-$F_1$ against a paragraph-length gold answer rewards
fluent medical prose whether or not it addresses the question, so the metric
cannot separate a good answer from a well-phrased refusal, and supports
\emph{relative} comparison between policies only, never absolute quality
claims. Open-ended conclusions are therefore weakly supported in
\emph{either} direction, which is a measurement limitation rather than a null
result, which a sharper metric (entailment against the gold claim, or human
adjudication) would be required to settle. The problem compounds with a
corpus--benchmark mismatch: PubMedQA asks about the findings of specific
studies, which the textbook-and-abstract corpus largely does not contain, so
retrieval raises the refusal rate to \retPoneOpen{} under always-retrieve
and the open-ended arm partly measures corpus coverage rather than gating.

\paragraph{Entropy gate omits the draft's first token.} The entropy gate
slices the logit buffer one row late, so $\bar{H}$ averages the draft's
second and subsequent tokens rather than all of them. The bias is
deterministic and identical across every question and both models, and
$\tau$ was calibrated on the same shifted quantity, so internal validity and
cross-model comparability are unaffected; what narrows is the definitional
claim. It was found too late to correct without re-running calibration and
evaluation on both models. See
Section~\ref{sec:limitation-entropy-offset} for the measurement and the
reasoning.

\paragraph{Hardware tiers not physically benchmarked.} The tier-adaptive
degradation mechanism is implemented and unit-tested, but only the medium
tier was actually run. The original project plan promised physical
benchmarking on three device classes; what was delivered is the mechanism
plus single-tier measurements, and this dissertation does not imply
otherwise.

\paragraph{Benchmark narrowness.} MCQ evaluation uses one MIRAGE subset
(MMLU-Med); open-ended evaluation uses PubMedQA only, after the originally
planned dataset proved unreachable. MIRAGE is entirely multiple-choice,
which is precisely why the open-ended complement matters, but one dataset
per question type limits generality.

\paragraph{Energy measurement.} Energy figures are \texttt{codecarbon}
software estimates, not metered hardware power measurements.

\paragraph{Risk tagger.} The clinical-risk tagger is a keyword rule set,
chosen deliberately for auditability, since a learned classifier would
obscure the basis of a safety-relevant label and no training data exists,
but it has recall gaps on implicitly high-risk questions. That gap is partly
\emph{why} the high-risk stratum is only $n=6$: the tagger finds explicit
dosing/interaction cues, not subtle criticality.

\paragraph{Stipulated safety bars.} The 70/80/90\% bars are stipulated, not
derived from clinical evidence. Different bars would move the envelope's
edges, though at current accuracy no defensible bar would be met.

\subsection{Off-by-one in the entropy gate's logit slice}
\label{sec:limitation-entropy-offset}

The entropy gate reads per-token logit distributions from
\texttt{llama-cpp-python}'s \texttt{\_scores} buffer, slicing generated-token
rows as \texttt{\_scores[n\_prompt : n\_prompt + max\_tokens]} where
\texttt{n\_prompt} is \texttt{len(tokenize(prompt))}. That start index is one
row too high: \texttt{tokenize()} prepends a beginning-of-sequence token the
completion's context does not contain, so the prompt occupies
$n_{\text{prompt}} - 1$ rows. The discrepancy was confirmed directly, since
under greedy decoding the $\arg\max$ of each logit row must be the token that
row produced, and for a 48-token draft the slice beginning at
$n_{\text{prompt}} - 1$ reproduced all 48 generated tokens while the slice
beginning at $n_{\text{prompt}}$ reproduced none. The practical consequence
is that $\bar{H}$ averages generated tokens $2 \ldots N$ rather than
$1 \ldots N$. On a representative clinical question the omitted first token
was the draft's highest-entropy one ($H = 3.28$ nats), and excluding it
lowered the gate signal from $0.763$ to $0.709$ against $\tau = 0.70$.

This does not invalidate the reported results. The bias is deterministic and
identical across every question and every model, so it is a consistent
redefinition of the signal rather than a source of noise; $\tau$ was
calibrated on the same shifted quantity, so the retrieval budget it produces
is the budget actually delivered; and the budget-matching procedure inherits
the identical bias at all three scales, keeping the cross-model comparison
like-for-like. What narrows is the definitional claim: the quantity this work
calls mean draft entropy is the mean over the draft's second and subsequent
tokens. Correcting it was judged the greater risk at the point of discovery,
since re-slicing would shift $\bar{H}$ upward on every question, move an
unknown number of near-threshold decisions across $\tau$, and require full
recalibration and re-evaluation on all three models. A corrected slice is the
first item of future work and is a single-line change once there is time to
re-run the calibration behind it.

\section{Implications for Offline Clinical Deployment}
\label{sec:disc-deployment}

What should a practitioner take from this work?

\paragraph{Do not deploy at this scale.} The unambiguous first implication:
a quantised 3B model, gated or not, is \gapToLowBar{}pp short of even the
lowest stipulated bar, and its safety envelope is empty. Scaling improves
but does not resolve this: the 7B and 14B clear the low-risk bar only, the
medium-risk bar remains out of reach at every scale tested, and the
capability curve is flattening as it approaches it
(Section~\ref{sec:scaling-bars}). This deployment class, meaning fully
offline, quantised, commodity-hardware clinical QA, is not fit for clinical
decision support beyond the lowest-risk stratum, and the honest output of
the evaluation is the measured distance rather than a product. Publishing
that distance is the responsible finding.

\paragraph{Where the method is useful today.} The gating approach earns its
keep wherever retrieval is the expensive or risky resource and the model's
self-knowledge is exploitable: cost-aware routing in on-premises deployments
under data-sovereignty constraints, edge devices where a metered or slow
retrieval backend dominates cost, and any pipeline where the operator wants
an auditable, per-query record of \emph{why} the system did or did not
consult external evidence. The gate's decision payload of signals, votes and
thresholds is logged per query, which makes the retrieval decision
reviewable after the fact in a way end-to-end trained routers are not.

\paragraph{Interpretability as the deployable asset.} The per-token entropy
attribution (Section~\ref{sec:eval-attribution}) points at what may matter
more than a small accuracy delta in regulated settings: surfacing
\emph{where} the model was uncertain. The honest caveat travels with it:
on this model, raw entropy highlights phrasing rather than clinical content,
so attribution is a promising direction with a demonstrated limitation, not
a finished feature. In a domain where automation bias is a documented harm
(Chapter~\ref{ch:lsep}), an uncertainty display that clinicians can audit is
plausibly worth more than two points of accuracy; building and validating
one is future work.
